% Algorithm text only; no document preamble or repeated classification derivation.
% Based on the algorithm paragraphs of the supplied PRL_SSG_Type_sup.pdf.
% Use the four blocks below for the enumeration paragraph in II.D,
% II.E.5, II.F.5, and II.G, respectively; adapt heading levels on insertion.
% Equation numbers cited in the prose refer to the original SM.
% Added displays are unnumbered and do not change its equation numbering.
% Required packages: amsmath, amssymb, bm.
\documentclass{article}
\usepackage{amsmath,amssymb,bm}

\begin{document}
\subsubsection{Algorithm for Type 0}

To enumerate all inequivalent $\rho_P$, we fix a normal subgroup
$H\lhd P$ and use $\rho_P(P)\simeq P/H$, as in Eq.~(12).
The enumeration proceeds as follows:
\begin{enumerate}
\renewcommand{\labelenumi}{(\roman{enumi})}
\item Enumerate all normal subgroups $H$ of $P$.
\item For each $H$, construct the quotient multiplication table and
determine the rotation groups $S_i\subset SO(3)$ abstractly isomorphic
to $P/H$.
\item For each allowed $S_i$, choose an isomorphism $P/H\to S_i$.
\item Lift this isomorphism to a homomorphism $\rho_P:P\to SO(3)$.
\item Merge spin-conjugate homomorphisms, combine the resulting
$\rho_P$ with the grading $\chi$, and impose spatial equivalence on
the complete representation $\varrho=\chi\rho$.
\end{enumerate}

These steps involve only finite groups. All normal subgroups are
obtained by starting from $\{e\}$ and repeatedly adjoining the normal
closure of an additional element. The possible images are $C_n$
($n=1,2,3,4,6$), $D_n$ ($n=2,3,4,6$), $T$, and $O$.
To determine an isomorphism, choose generators of $P/H$ and enumerate
their images in $S_i$, requiring the same element orders and generation
of the whole target. Extend each assignment by multiplication. Discard
an assignment if two words representing the same quotient element
have different images. A surviving map $f$ is accepted only if it is
bijective and satisfies
\[
 f(qq')=f(q)f(q')\qquad\text{for all }q,q'\in P/H.
\]
Thus group orders and element orders provide preliminary exclusions,
while the multiplication law establishes the isomorphism.

For the rotation groups in this list, every automorphism of $S_i$
is induced by its $SO(3)$ normalizer. For cyclic groups the only
generator changes are identity and inversion; for the dihedral groups
the generator changes are realized by rotations of the principal and
transverse axes. The normalizer of $D_2$ permutes its three axes,
the normalizer $O$ realizes all automorphisms of $T$, and the
automorphisms of $O$ are inner. Consequently, one isomorphism for
each allowed quotient suffices after spin conjugation. Different
kernels are retained until the spatial quotient is taken.

\subsubsection{Algorithm for Type 1}

The procedure to enumerate Type 1 SSGs with parent space group $G$ is:
\begin{enumerate}
\renewcommand{\labelenumi}{(\roman{enumi})}
\item Enumerate all normal subgroups $P_0\lhd P$ of index two;
each defines a grading $\eta:P\to\{\pm1\}$. Include the trivial
grading as well.
\item For each $\eta$, determine the wave vectors satisfying the
covariant condition (19). The translation holonomy is
$\rho_{\mathrm T}(t)=e^{-i(\bm k\cdot t)J_z}$.
\item Solve the obstructed 1-cocycle equation (20). The existence of
a solution tests $[\nu_{\bm k}]=1$, and its solutions determine the
section matrices through Eq.~(17). Construct
$\rho(t s(p))=\rho_{\mathrm T}(t)\rho(s(p))$.
\item Merge spin-conjugate homomorphisms and identify continuous
families by the complete wave-vector and section-phase solutions,
as described below.
\item Combine the resulting $\rho$ with $\chi$ and impose spatial
equivalence on $\varrho=\chi\rho$.
\end{enumerate}

The second and third steps can be carried out together by integer
linear algebra. In a primitive basis $\{\bm a_j\}$, write
$\bm k=\sum_j q_j\bm b_j$, with
$\bm a_i\cdot\bm b_j=2\pi\delta_{ij}$, and
$t_{p_1,p_2}=\sum_j c_j(p_1,p_2)\bm a_j$. Let $M_p$ be the
integer matrix of $R_p$ in this basis and set
$x_p=\theta(s(p))/(2\pi)$. The required equations are
\[
 (M_p^T-\eta_p I)\bm q=0\pmod1,\qquad
 x_{p_1}+\eta_{p_1}x_{p_2}-x_{p_1p_2}
       -\bm q\cdot\bm c(p_1,p_2)=0\pmod1.
\]
A common spin rotation changes $x_p$ by $(1-\eta_p)v$.
For nontrivial $\eta$, choose $r$ with $\eta_r=-1$ and fix this
gauge by $x_r=0$. For trivial $\eta$ there is no continuous gauge
direction. The remaining conjugation by $R_x(\pi)$ simultaneously
sends $(\bm q,x)$ to $(-\bm q,-x)$.

Stacking the equations, including the gauge condition, gives an
integer congruence $Az=b\pmod1$. For a specified wave vector,
$z=x$ and $b$ contains its extension phases; for the joint problem,
$z=(\bm q,x)$ and $b=0$. Compute the Smith form $UAV=D$, with
$U,V$ integer unimodular and positive nonzero diagonal entries
$d_1,\ldots,d_r$. Writing $z=Vy$, all solutions are
\[
 y_i=\frac{(Ub)_i+m_i}{d_i}\pmod1,
 \quad m_i=0,\ldots,d_i-1\quad (i\le r),
 \qquad (Ub)_i=0\pmod1\quad(i>r).
\]
The last conditions refer to the zero rows of $D$ and decide
solvability. Coordinates corresponding to zero columns are free in
$\mathbb R/\mathbb Z$. Thus an inconsistent system detects the
obstruction, while a consistent system gives every solution, including
all torsion branches. At fixed $\bm k$, the gauge-inequivalent solutions
form a torsor for $H^{1,\eta}(P,U(1))$: after one particular solution
is chosen, all others are obtained by adding homogeneous solutions.

For the joint system, retain the resulting disjoint affine tori
\[
 X_\eta=\coprod_a C_a,\qquad
 C_a=z_a+B_a(\mathbb R/\mathbb Z)^{d_a}.
\]
The offsets $z_a$ are rational and the matrices $B_a$ are integer.
Intersect these components with $2\bm q=0$ modulo one. The
$\bm q=0$ representations are counted in Type 0, and nonzero-TRIM
representations are retained separately. The remaining part of each
positive-dimensional component defines a general family; isolated
non-TRIM solutions remain discrete entries. Keeping the joint
solutions is essential because a reciprocal winding can permute
section-phase branches.

\subsubsection{Algorithm for Type 2}

For a given parent space group $G$, all Type 2 SSGs are enumerated
as follows:
\begin{enumerate}
\renewcommand{\labelenumi}{(\roman{enumi})}
\item Determine every nonzero $t_0\in L/2L$ invariant under $P$
modulo $2L$. For each allowed $t_0$, construct $\rho_{\mathrm T}$
from its annihilator triple $\{\bm k_i\}$.
\item Determine $\sigma:P\to S_3$ from
$R_p\bm k_i=\bm k_{\sigma_p(i)}$ modulo the reciprocal lattice,
and construct $S(p)$.
\item Solve Eqs.~(37) and (38) for $D(s(p))$, and construct
$\rho(t s(p))=\rho_{\mathrm T}(t)D(s(p))S(p)$.
\item Merge homomorphisms related by a simultaneous spin
conjugation of their translation and section matrices.
\item Combine the resulting $\rho$ with $\chi$ and impose spatial
equivalence on the complete $\varrho$.
\end{enumerate}

In the second step, let $P_\sigma\bm e_i=\bm e_{\sigma(i)}$ and
choose $S_\sigma=\operatorname{sgn}(\sigma)P_\sigma$.
Then $S_\sigma\in SO(3)$ and $S_\sigma S_\tau=S_{\sigma\tau}$,
so the lift introduces no additional factor system.
For the third step, put
$d_{p,i}=\theta_i^s(p)/\pi$ and
$b_{p_1p_2,i}=\bm k_i\cdot t_{p_1,p_2}/\pi$ modulo two.
Equations (37) and (38) become a finite linear system over
$\mathbb F_2$:
\[
 d_{p_1,i}+d_{p_2,\sigma_{p_1}^{-1}(i)}-d_{p_1p_2,i}
       =b_{p_1p_2,i},\qquad
 \sum_i d_{p,i}=0.
\]
Gaussian elimination either detects an inconsistent row or gives a
particular solution $d^{(0)}$ and a basis $h_1,\ldots,h_r$ of the
homogeneous kernel. Every solution is then
\[
 d=d^{(0)}+\sum_{j=1}^r\epsilon_jh_j,
 \qquad \epsilon_j\in\mathbb F_2.
\]
Conjugation by a fixed element of $V_4$ adds a coboundary. Quotienting
by these coboundaries gives the torsor for $H^{1,\pi}(P,V_4)$;
axis permutations must also be included. Equivalently, as done here,
construct all affine solutions and compare their full sets of matrices
under one common signed permutation. This removes both redundancies
and does not identify matrices using independent basis changes.

\subsection{Equivalence relations}

Three kinds of equivalence must be distinguished.

\textit{(1) Spin-space conjugation.}
Two representations are equivalent if one common spin transformation
conjugates all their matrices. For $O(3)$ representations, conjugation
by $Q$ and $-Q$ is identical, so $SO(3)$ transformations suffice.
The Type 0 construction removes the isomorphism choices described
above. Type 1 requires the continuous phase gauge and simultaneous
reversal of wave vector and phases. Type 2 requires simultaneous
signed permutations of the translation and section matrices.

\textit{(2) Automorphisms of $G$.}
Following the chirality-preserving convention, impose
$\varrho\sim\varrho\circ\alpha$ for affine coordinate changes
$\{A\mid\bm o\}$ normalizing $G$, with $A\in SL(3,\mathbb Z)$
in a primitive basis. They satisfy
\[
 A M_p A^{-1}=M_{\sigma(p)},\qquad
 \bm u_p=A\bm\tau_p+(I-M_{\sigma(p)})\bm o
                      -\bm\tau_{\sigma(p)}\in\mathbb Z^3.
\]
Obtain generators of the integral point-group normalizer, restrict
them to the actual primitive lattice, and take the determinant-one
subgroup. Triclinic and monoclinic groups require generators of the
appropriate integral linear groups; they cannot be treated by a
bounded search over coordinate matrices. For the remaining crystal
systems the linear normalizer is finite and can be enumerated from
its invariant axes or invariant positive-definite metric.

One must next retain the stabilizer of the nonsymmorphic extension,
rather than only those individual generators that already preserve
it. Stack the section translations into $\tau$ and the matrices
$I-M_p$ into $B_0$. If the rows of $C_0$ span the full integer
annihilator of $B_0$, then $C_0\tau$ modulo one represents the
extension data modulo origin shifts. Its orbit under the linear
normalizer is finite because the section translations have a common
finite denominator. Traverse this orbit. If $t_x$ carries the initial
class to $x$ and a generator $s$ carries $x$ to $y$, the elements
$t_y^{-1}s t_x$ generate its stabilizer. For each such element $h$,
solve $B_0\bm o=\tau-\tau^h$ modulo one by Smith reduction to obtain
an affine lift, where $\tau^h_{\sigma(p)}=h\bm\tau_p$.
Include also the discrete components of $B_0\bm o=0$ modulo one.
Continuous origin shifts in its kernel centralize $G$; integer origin
shifts induce inner automorphisms and hence spin conjugations.

An affine generator pulls back the Type 1 data according to
\[
 \bm q'=A^T\bm q,\qquad
 \eta'_p=\eta_{\sigma(p)},\qquad
 x'_p=x_{\sigma(p)}+\bm q\cdot\bm u_p\pmod1.
\]
Restore the spin gauge and compare the complete affine components,
including their phases. The same automorphism must act on $\chi$;
separate spatial quotients of $\rho$ and $\chi$ are not sufficient.
For discrete representations, apply the same pullback to all
translation and section matrices.

\textit{(3) Continuous families.}
A continuous family is determined by continuation of the complete
general representation branch, with its discrete constituent data
fixed. In the joint solution above, each positive-dimensional $C_a$
carries one general-family label after its special representations
have been recorded separately. Equal $\eta$ or a common wave-vector
domain alone does not establish equivalence between distinct
components. At fixed wave vector an infinitesimal phase variation
satisfies $z_{p_1}+\eta_{p_1}z_{p_2}-z_{p_1p_2}=0$. Averaging over
$p_2$ gives
\[
 z_p=(1-\eta_p)\bar z,\qquad
 \bar z=|P|^{-1}\sum_{p\in P}z_p.
\]
Thus the only continuous phase freedom is the spin gauge already
removed. The projection of $C_a$ onto its wave-vector image $K_a$
is consequently a connected finite covering, retaining the phase
changes under reciprocal-period winding.

Removing special wave vectors may disconnect the general part of
$C_a$ without producing new continuation families. Such pieces are
identified by continuation of the same branch formula, including
through a zero wave vector, or by its reciprocal monodromy. The
special representations themselves retain their separate labels.
This prescription does not permit switching between unrelated
branches merely because their limiting matrices coincide.

With the same chirality convention, this deduplication is equivalent
to the representation-classification prescription of Song and
co-workers.\footnote{Z. Xiao, J. Zhao, Y. Li, R. Shindou, and
Z.-D. Song, \textit{Spin Space Groups: Full Classification and
Applications}, Phys. Rev. X \textbf{14}, 031037 (2024), Sec.~II D
and Appendix~B.}
Spin conjugation and spatial pullback are the same operations in
both descriptions. A general Type 1 block gives the real form of a
complex character, or of its induction from $\ker\eta$; conversely,
these real two-dimensional blocks give the joint equations above.
The residual real character must also be retained: if
$\varrho=\beta\oplus\nu$, then
$\chi=(\det\beta)\nu$ and $D=\chi\beta$ give
$\varrho=\chi\operatorname{diag}(D,\det D)$.
This invertible change of description includes the translation-character
shift between the wave vectors of $\beta$ and $D$.

Song's continuous identifications are generated by motion along a
fixed branch, exchange of conjugate wave vectors, and reciprocal
windings. They lift to paths in $C_a$, together with the spin
identification $(\bm q,x)\sim(-\bm q,-x)$. Conversely, the monodromy
of the connected covering $C_a\to K_a$ is transitive on its phase
sheets and is generated by a primitive set of reciprocal periods
of $K_a$. It therefore yields precisely these branch identifications,
with special points used only for continuation of the same branch.
Finally, apply the spatial and spin generators to the discrete
representations and continuous families, merging their classes until
no new identifications occur. This transitive closure implements the
same three equivalences; equality of family labels does not assert
spin conjugacy at arbitrary fixed parameter values.

\end{document}